% Nouveau EGA V (cote 25): archive cote 25, site page 16 (PDF page 17).
% typescript. Draft compared: none.
% AI-transcribed by Claude Opus 5.5 (claude-opus-5-5, adaptive thinking, effort high), Message Batches API, batch ega25b (msgbatch_01A7n5YuXubUPNK1zsepi8RG), 2026-11-13; original AI transcription, subsequently checked against both supplied scan parts; see PAGES_CHECKED.tsv and CHECK_DECISIONS.tsv.
\scanpage{25/16}
\mspage{\hand{13}}

 Soit $K/k$ le corps des fonctions de $X$, et soit $n=\dim P$;
introduisant des coordonnées  affines $T_1,..,T_n$ dans $P$ (en choisiss
ant un hyperplan à l'infini $H_\hand{\infty}$ tel que $f(X)$ ne soit pas contenu dans
celui-ci) et $S_1,...S_n$ dans $\overset{\hand{\vee}}{P}$ , on voit que le corps des fonctions $L$
de $Y_{\hand{\eta}}$  s'identifie au corps des fractions de l'anneau intègre
$K[S_1,..,S_n]/(\hand{\sum} t_iS_i - 1)$ , où les $t_i \hand{\in} K$ sont les images des $T_i$ à
l'adde de $f:X\hand{\to} P$ . Comme $\dim \hand{\overline{\textcolor{black}{f(X)}}} \hand{>} 0$ , les $t_i$ ne sont pas tous
algébriques sur $k$, à fortiori ils ne sont pas tous nuls, supposons
p. ex. $t_n \hand{\neq} 0$. On constate alors aussitôt qu'on a $L = K(S_1,..,S_{n-1})$ ,
(extension transcendante pure), $S_n\hand{\in} L$ étant donné par l'équation
\[ \hand{\sum} t_iS_i - 1 = 0 \]
en fonction des $S_i$ $(1\hand{\leq} i\hand{\leq} n-1)$ et des $t_i$ $(1\hand{\leq} i\hand{\leq} n)$. D'autre part,
$k' = k(\hand{\eta})$ s'identifie à $k(S_1,..,S_n)$, et l'inclusion canonique
$k'\hand{\longrightarrow} L$ s'obtient en envoyant $S_i$ sur $S_i$ , i.e. $k'$ en tant que sous-%
extension de $L$ est la sous-extension engendrée par les $S_i$ $(1\hand{\leq} i\hand{\leq} n)$, ou
ce qui revient évidemment au même, par les $S_i$ $(1\hand{\leq} i\hand{\leq} n-1)$ et par
$S_n = a_o + a_1S_1 + \ldots + a_{n-1}S_{n-1}$ , où $a_o = t_n^{-1}$ , $a_i =-t_it_n^{-1}$ pour
$1\hand{\leq} i\hand{\leq} n-1$. Notons que le corps engendré par les $a_i$ ou par les $t_i$ est évidem
ment le même, leur degré de transcendance commun
n'est autre que $\dim f(X)$.

\placedmargin{16}
